\documentclass[exercice]{thedocument}%
\begin{document}%
    \>S1=F;\>S2=G;\>S3=H;\>S4=I;\>S5=J;%
    \enonce{0046}%
    \begin{center}
        \begin{tikzpicture}
            \coordinate[label=\<S1;](S1)at(-2,0);
            \coordinate[label=\<S3;](S3)at(2.5,1);
            \path(S1)--(S3)coordinate[pos=-0.2](d0);
            \path(S1)--(S3)coordinate[pos=1.2](d1);
            \coordinate[label=above right:\<S2;](S2)at(-1,1.5);
            \coordinate[label=below right:\<S4;](S4)at(2,-1);
            \draw(d0)--(d1);
            \foreach\ii in{1,2,3,4}{%
                \node at(S\ii){$\times$};
            }
        \end{tikzpicture}
    \end{center}
\end{document}